\documentclass[12pt,a4paper]{article}
\usepackage[utf8]{inputenc}
\usepackage[T1]{fontenc}
\usepackage{amsmath,amssymb}
\usepackage{graphicx}

\begin{document}

\section*{Zadání}

Nechť $A, B, C$ jsou matice typu $2\times 2$:
\[
A=\begin{pmatrix} a_{11} & a_{12}\\ a_{21} & a_{22}\end{pmatrix},\quad
B=\begin{pmatrix} b_{11} & b_{12}\\ b_{21} & b_{22}\end{pmatrix},\quad
C=\begin{pmatrix} c_{11} & c_{12}\\ c_{21} & c_{22}\end{pmatrix}.
\]
Určete podmínky, které musí tyto matice splňovat, aby existovala nenulová matice
$X=\begin{pmatrix} x_{11} & x_{12}\\ x_{21} & x_{22}\end{pmatrix}$ typu $2\times 2$, pro kterou platí
\[
X = B\left(X - CA^{T}\right) + \left(AC^{T}B^{T}\right)^{T}.
\]

\section*{Řešení}

Rovnici začneme řešit za využití pravidel pro počítání s maticemi:
\begin{enumerate}
  \item $AB \neq BA$ (obecně),
  \item $(AB)C = A(BC)$,
  \item $(AB)^{T} = B^{T}A^{T}$,
  \item $X\in\mathbb{R}^{n\times n},\ AX=E_n \ \wedge\ XA=E_n \ \Rightarrow\ X=A^{-1}$.
\end{enumerate}

\medskip
Vyjdeme z rovnice
\[
X = B\left(X - CA^{T}\right) + \left(AC^{T}B^{T}\right)^{T}.
\]
Víme následující:
\[
B\left(X - CA^{T}\right) \iff BX - BCA^{T},
\qquad
C^{T}B^{T} \iff (BC)^{T}.
\]
Dosadíme tedy do rovnice:
\[
X = BX - BCA^{T} + \left[A(BC)^{T}\right]^{T}.
\]
Dále označíme
\[
(BC)^{T} = D
\]
a po dosazení získáme
\[
X = BX - BCA^{T} + (AD)^{T},
\]
neboli
\[
X = BX - BCA^{T} + D^{T}A^{T}.
\]
Po zpětném dosazení $(BC)^{T}=D$ získáme
\[
X = BX - BCA^{T} + \left[(BC)^{T}\right]^{T}A^{T},
\]
přičemž z definice transpozice
\[
A^{T} =
\begin{pmatrix}
a_{11} & a_{12} & \cdots & a_{1n}\\
\vdots & \vdots & \ddots & \vdots\\
a_{m1} & a_{m2} & \cdots & a_{mn}
\end{pmatrix}^{T}
=
\begin{pmatrix}
a_{11} & \cdots & a_{m1}\\
a_{12} & \cdots & a_{m2}\\
\vdots & \ddots & \vdots\\
a_{1n} & \cdots & a_{mn}
\end{pmatrix}.
\]
Z toho plyne, že opětovnou transpozicí již jednou transponované matice získáme opět původní matici, neboli
\[
\left(M^{T}\right)^{T} = M,
\]
a tudíž (pro $M = BC$)
\[
X = BX - BCA^{T} + \left[(BC)^{T}\right]^{T}A^{T},
\]
\[
X = BX - BCA^{T} + BCA^{T}.
\]
Z čehož plyne, že
\[
X = BX,
\]
a tudíž matice $C, A$ mohou být zcela libovolné.

\section*{Podmínka pro matici $B$}

Pro zachování rovnosti $X = BX \ \wedge\ X \neq 0$ musí nastat jeden z případů:
\begin{enumerate}
  \item[I.] $B = E$, protože platí $E\cdot X = X$, kde X může být libovolná nenulová 2x2
  \item[II.] $B = \begin{pmatrix} 0 & 1\\ 1 & 0\end{pmatrix}$, pak
  \[
  BX =
  \begin{pmatrix} 0 & 1\\ 1 & 0\end{pmatrix}
  \begin{pmatrix} x_{11} & x_{12}\\ x_{21} & x_{22}\end{pmatrix}
  =
  \begin{pmatrix} x_{21} & x_{22}\\ x_{11} & x_{12}\end{pmatrix},
  \]
  a rovnost $X = BX$ tedy platí pro každou nenulovou matici $X$, jejíž prvky jsou všechny stejné,
  tedy $x_{11} = x_{12} = x_{21} = x_{22}$, např. $X=\begin{pmatrix} 1 & 1\\ 1 & 1\end{pmatrix}$.
  \item[III.] $B = X$, pak z rovnosti $BX = X$ dostaneme $X\cdot X = X$, tedy
  \[
  X^{2} = X,
  \]
  a matice $B = X$ musí být nenulová matice splňující $B^{2} = B$ ,
  např. $B = X = \begin{pmatrix} 1 & 0\\ 0 & 0\end{pmatrix}$.

  Najdeme všechny matice $X$, pro které platí $X^{2}=X$.
  
  \item Vynásobíme obecnou matici sama se sebou:
  \[
  \begin{pmatrix} x_{11} & x_{12}\\ x_{21} & x_{22}\end{pmatrix}
  \begin{pmatrix} x_{11} & x_{12}\\ x_{21} & x_{22}\end{pmatrix}
  =
  \begin{pmatrix}
  x_{11}^{2}+x_{12}x_{21} & x_{11}x_{12}+x_{12}x_{22}\\
  x_{21}x_{11}+x_{22}x_{21} & x_{21}x_{12}+x_{22}^{2}
  \end{pmatrix}.
  \]
  Tato matice se musí rovnat matici $X$, takže dostaneme soustavu
  \begin{align*}
  x_{11}^{2}+x_{12}x_{21} &= x_{11},\\
  x_{11}x_{12}+x_{12}x_{22} &= x_{12},\\
  x_{11}x_{21}+x_{21}x_{22} &= x_{21},\\
  x_{12}x_{21}+x_{22}^{2} &= x_{22}.
  \end{align*}

  Z první rovnice plyne
  \begin{equation}
  x_{12}x_{21} = x_{11}-x_{11}^{2} = x_{11}(1-x_{11}).
  \end{equation}
  Ve druhé rovnici vytkneme $x_{12}$ a ve třetí $x_{21}$:
  \begin{align}
  x_{12}\,(x_{11}+x_{22}-1) &= 0,\\
  x_{21}\,(x_{11}+x_{22}-1) &= 0.
  \end{align}
  Pokud je alespoň jedno z čísel $x_{12}, x_{21}$ nenulové, musí tedy platit
  $x_{11}+x_{22}=1$, neboli
  \begin{equation}
  x_{22} = 1-x_{11},
  \end{equation}
  a z (1) dostaneme
  \begin{equation}
  x_{12}x_{21} = x_{11}(1-x_{11}).
  \end{equation}
  Čtvrtá rovnice je pak splněna automaticky, protože
  $x_{12}x_{21}=x_{22}-x_{22}^{2}=(1-x_{11})\,x_{11}$.

  Zbývá případ $x_{12}=x_{21}=0$. Soustava se zredukuje na $x_{11}^{2}=x_{11}$ a $x_{22}^{2}=x_{22}$,
  tedy $x_{11},x_{22}\in\{0,1\}$. Nenulová řešení jsou
  $\begin{pmatrix} 1 & 0\\ 0 & 0\end{pmatrix}$,
  $\begin{pmatrix} 0 & 0\\ 0 & 1\end{pmatrix}$ a $E$.
  První dvě už spadají do tvaru výše, třetí je $X=E$ (případ I.).

  \textbf{Výsledek:} nenulová řešení rovnice $X^{2}=X$ jsou
  \[
  X=\begin{pmatrix} x_{11} & x_{12}\\ x_{21} & 1-x_{11}\end{pmatrix},
  \qquad\text{kde } x_{12}x_{21}=x_{11}(1-x_{11}).
  \]
  
\end{enumerate}


\end{document}
